% !PRINTED 391
\weberpage{391}
Ist die Periode der Kettenbruchentwickelung von $\omega$
\weq{17}{[\alpha_0, \alpha_1 \ldots \alpha_{\nu-1}],}
so ist sie für $-1 : \omega'$:
\weq{18}{[\alpha_{\nu-1}, \alpha_{\nu-2} \ldots \alpha_0],}
also die umgekehrte.

Ist nun $\omega$ mit $\omega'$, also auch mit $-1 : \omega'$ äquivalent, so müssen
nach 6. diese beiden Perioden mit einander übereinstimmen, wenn
man bei einem geeigneten Gliede beginnt, oder die Periode von
$\omega$ muss \gesperrt{umkehrbar} sein.

Das besagt: Es müssen sich in der Kettenbruchentwickelung
für $\omega$ zwei Elemente so auswählen lassen, dass die Entwickelung
gleich lautet, wenn man sie von der einen dieser Zahlen nach
links oder von der anderen nach rechts fortschreitend liest.
In Zeichen: es muss für irgend ein passend bestimmtes $k$ und
für $i = 0, 1, 2 \ldots \nu - 1$
\[
\alpha_i = \alpha_{\nu+k-i}
\]
sein. Wenn die Periode umkehrbar ist, so ist auch umgekehrt
(nach 6.) $\omega$ mit $-1 : \omega'$, also auch mit $\omega'$ äquivalent. Daraus
also das Resultat:

8. \gesperrt{Zweiseitige Zahlen haben in ihrer Kettenbruchentwickelung
 eine umkehrbare Periode und umgekehrt
 sind Zahlen, die in der Kettenbruchentwickelung
 eine umkehrbare Periode haben,
zweiseitig.}

Schliesslich wollen wir noch bemerken, dass durch die Eigenschaft,
in einen periodischen Kettenbruch entwickelbar zu sein,
die reellen quadratischen Irrationalzahlen von allen anderen
Zahlen unterschieden sind. Denn ist $\omega$ eine in einen Kettenbruch
 entwickelte Zahl, und sind $\omega_n$, $\omega_m$ zwei Schlusszahlen
dieses Kettenbruches, so ist
\[
\omega = \frac{P_n\omega_n + P_{n-1}}{Q_n\omega_n + Q_{n-1}} = \frac{P_m\omega_m + P_{m-1}}{Q_m\omega_m + Q_{m-1}}.
\]

Wenn nun der Kettenbruch periodisch ist, gleichviel, ob
die Periodicität gleich von Anfang beginnt oder erst im Verlaufe
der Entwickelung, so ist für zwei verschiedene Werthe von $n$
und $m$, $\omega_n = \omega_m$, und man erhält also durch Elimination von
%CONT>>
